% !TEX root = ../../tree_notes.tex
% =====================================================================
%  1.4  Counting Trees                                       sec:c01-counting
%  Exports: def:c01-shape def:c01-labeled
%           thm:c01-catalan-shapes  thm:c01-cayley
%           lst:c01-catalan  lst:c01-cayley-pow  prob:c01-unique-bst
% =====================================================================
\section{Counting Trees}\label{sec:c01-counting}

\begin{topicintro}
"How many trees are there?" has two very different answers depending on what you are
counting: the number of \emph{shapes} (Catalan numbers) and the number of ways to put
\emph{distinct labels} on those shapes (Cayley's $n^{\,n-2}$). Both show up as DP counts and
as combinatorics problems, so it pays to know the two formulas and which question each answers.

\textbf{Prerequisites.} \secref{c01-kinds} (binary shapes) and basic binomial coefficients.

\textbf{After this section you can}
\begin{itemize}[nosep]
  \item count binary-tree shapes with the Catalan recurrence and closed form;
  \item state Cayley's formula and say what a labelled tree is;
  \item compute both modulo a prime without overflow;
  \item recognise a Catalan count hiding inside a DP problem (\chref{c10}).
\end{itemize}
\end{topicintro}

\subsection{Intuition}

An \term{ordered} binary tree distinguishes a left child from a right child, so a
three-node tree can be arranged in five different shapes -- one per way the single "extra"
node can hang off a two-node skeleton.

\begin{figure}[H]
\centering
\begin{tabular}{@{}c@{\hspace{2mm}}c@{\hspace{2mm}}c@{\hspace{2mm}}c@{\hspace{2mm}}c@{}}
\begin{forest} snaptree [1 [2 [3] [,phantom]] [,phantom]] \end{forest} &
\begin{forest} snaptree [1 [2 [,phantom] [3]] [,phantom]] \end{forest} &
\begin{forest} snaptree [1 [2] [3]] \end{forest} &
\begin{forest} snaptree [1 [,phantom] [2 [3] [,phantom]]] \end{forest} &
\begin{forest} snaptree [1 [,phantom] [2 [,phantom] [3]]] \end{forest} \\
{\scriptsize (a)} & {\scriptsize (b)} & {\scriptsize (c)} & {\scriptsize (d)} & {\scriptsize (e)}
\end{tabular}
\caption{All five shapes of a three-node binary tree, so $C_3 = 5$. They differ only in
which of the four slots the nodes occupy -- empty slots are drawn with \code{[,phantom]}.}\label{fig:c01-shapes3}
\end{figure}

\subsection{Definitions and Notation}

\begin{definition}{Shape vs labelled tree}{c01-shape}
The \term{shape} of a rooted (binary) tree is the tree with its values ignored -- only the
parent--child connections matter. A \term{labelled tree} on $n$ vertices keeps the shape
\emph{and} assigns the distinct names $1,2,\dots,n$ to the nodes. Two trees are equal as
shapes iff an isomorphism maps one to the other preserving children; they are equal as
labelled trees only if the labels match too.
\end{definition}

\begin{definition}{Counted objects}{c01-labeled}
We count three related quantities for $n$ nodes: the number of \emph{binary-tree shapes}
(ordered, values ignored); the number of \emph{binary trees labelled} $1..n$ (a shape times
the ways to fill it); and the number of \emph{labelled trees} whose edges are chosen freely
(not restricted to be binary, and undirected). Each has its own formula below.
\end{definition}

\subsection{Shapes: the Catalan Numbers}

\begin{theorem}{Binary Tree Shapes are Catalan}{c01-catalan-shapes}
The number of distinct shapes of a binary tree with $n$ nodes is the Catalan number
\[
  C_n \;=\; \frac{1}{n+1}\binom{2n}{n} \;=\; 1,\,1,\,2,\,5,\,14,\,42,\,132,\dots
  \quad (n = 0,1,2,\dots)
\]
equivalently $C_0 = 1$ and $C_n = \sum_{i=0}^{n-1} C_i\,C_{n-1-i}$.
\end{theorem}
\begin{proofbox}
Root-splitting counts the recurrence. A tree with $n\ge 1$ nodes spends one node on the
root; the remaining $n-1$ split into $i$ nodes in the left subtree and $n-1-i$ in the
right, for some $i = 0,\dots,n-1$. Left and right shapes are chosen independently, giving
$C_i \cdot C_{n-1-i}$ trees for that $i$, and summing over $i$ yields
$C_n = \sum_{i} C_i C_{n-1-i}$ with $C_0 = 1$ (the empty shape). This is exactly the
Catalan recurrence, whose closed form is $\frac{1}{n+1}\binom{2n}{n}$ (obtainable by
generating functions or by a bijection with balanced parenthesis strings).
\end{proofbox}

\begin{figure}[H]
\centering
\begin{tikzpicture}[diagram]
  \node[troot] (r) at (0,0) {root};
  \node[tnode] (l) at (-2.2,-1.6) {\shortstack{left\\$i$ nodes}};
  \node[tnode] (rr) at (2.2,-1.6) {\shortstack{right\\$n{-}1{-}i$ nodes}};
  \draw[tedge] (r) -- (l); \draw[tedge] (r) -- (rr);
  \node[annot] at (0,0.9) {$1$ node used by the root};
  \node[annot] at (0,-2.7) {choose $i = 0 \ldots n-1$;  shapes $= C_i \cdot C_{n-1-i}$};
\end{tikzpicture}
\caption{The Catalan recurrence: the root splits the other $n-1$ nodes between two
independent subtrees, and we sum the products over every split.}\label{fig:c01-catalan-split}
\end{figure}

\begin{notebox}[The number of labelings for a fixed shape]
Given a binary shape of $n$ nodes, filling it with labels $1..n$ so that it stays a
\emph{BST} is forced (in-order must be sorted), which is why counting BSTs $=$ counting
shapes $= C_n$. Filling it with labels with no BST rule gives $n!$ labelled trees on that
shape.
\end{notebox}

\subsection{Labelled Trees: Cayley's Formula}

\begin{theorem}{Cayley's Formula}{c01-cayley}
The number of distinct trees on $n$ \emph{labelled} vertices $1,\dots,n$ (edges arbitrary,
not restricted to binary or rooted) is $n^{\,n-2}$. Rooted at any vertex, this becomes
$n^{\,n-1}$ labelled rooted trees.
\end{theorem}
\begin{proofbox}[Proof idea (Pr\"ufer sequences)]
There is a bijection between labelled trees on $n$ vertices and sequences ("Pr\"ufer codes")
of length $n-2$ whose entries are any of the $n$ labels. Since each of the $n-2$ positions
holds one of $n$ choices freely, there are $n^{\,n-2}$ such sequences, hence that many
trees. The bijection deletes the smallest leaf repeatedly and records its neighbour;
reversing the process rebuilds the tree from the code. The construction and its inverse
are built in full in \chref{c17}.
\end{proofbox}
\begin{insightbox}
$n=1$: one tree ($1^{-1}$ read as $1$). $n=2$: $2^0=1$ edge. $n=3$: $3^1=3$. $n=4$:
$4^2=16$. Growing far slower than $n!$ -- most permutations of a shape collapse to the
same unrooted tree.
\end{insightbox}

\usedin{Chapter 10 and Chapter 20 (DP that counts trees), and Chapter 17 (Pr\"ufer codes)}
{Catalan-style root-splitting recurrences reappear directly in tree-counting DPs, and the
Pr\"ufer bijection behind Cayley's formula -- including building a tree from its code --
is constructed in Chapter 17.}

\subsection{Computing the Counts}

\begin{lstlisting}[caption={Catalan numbers by DP, $C_0 \dots C_n$ (grows fast -- use big integers past $n \approx 30$)}, label={lst:c01-catalan}]
#include <bits/stdc++.h>
using namespace std;

// C[0..n], the number of binary-tree shapes with k nodes is C[k].
vector<long long> catalan(int n) {
    vector<long long> C(n + 1, 0);
    C[0] = 1;
    for (int k = 1; k <= n; ++k)
        for (int i = 0; i < k; ++i)
            C[k] += C[i] * C[k - 1 - i];     // root split: i left, k-1-i right
    return C;
}
\end{lstlisting}

\begin{lstlisting}[caption={Cayley modulo a prime: $n^{\,n-2} \bmod m$ by fast exponentiation}, label={lst:c01-cayley-pow}]
long long modpow(long long b, long long e, long long m) {
    long long r = 1 % m;
    b %= m;
    for (; e > 0; e >>= 1, b = b * b % m)
        if (e & 1) r = r * b % m;
    return r;
}
// number of labelled trees on n vertices modulo m
long long labelledTrees(long long n, long long m) {
    return n <= 1 ? 1 % m : modpow(n, n - 2, m);   // n^(n-2) mod m
}
\end{lstlisting}

\begin{complexitybox}
The Catalan DP is $O(n^2)$ time, $O(n)$ memory; use $\bigO{n}$ only if you switch to the
closed form with modular inverses. Cayley with binary exponentiation is $O(\log n)$ after
the $n-2$ exponent is set. Both counts overflow \code{long long} almost immediately
($C_{35}$ exceeds $2^{63}$), so contests ask for the answer modulo a prime -- then reduce
every multiply.
\end{complexitybox}

\subsection{Worked Examples}

\begin{example}{Small values}{c01-counting-ex}
\figref{c01-shapes3} shows $C_3 = 5$. Extending the root-split,
$C_4 = C_0C_3 + C_1C_2 + C_2C_1 + C_3C_0 = 5 + 2 + 2 + 5 = 14$, and $C_5 = 42$.
For labelled trees, Cayley gives $n = 4 \Rightarrow 4^{2} = 16$ and $n = 5 \Rightarrow 5^{3} = 125$.
\end{example}

\begin{tracebox}
\lstname\ of \coderef{c01-catalan} builds the array $C = [1, 1, 2, 5, 14, 42, 132]$ for
$n = 6$: each new entry is the running sum of products from \figref{c01-catalan-split},
matching \thmref{c01-catalan-shapes} line for line.
\end{tracebox}

\begin{center}
\renewcommand{\arraystretch}{1.1}
\begin{tabular}{@{}c|ccccccc@{}}
\textbf{$n$}   & 0 & 1 & 2 & 3 & 4 & 5 & 6 \\ \hline
$C_n$ (shapes) & 1 & 1 & 2 & 5 & 14 & 42 & 132 \\
labelled trees $n^{n-2}$ & -- & 1 & 1 & 3 & 16 & 125 & 1296 \\
\end{tabular}
\end{center}

\subsection{Common Pitfalls}

\begin{warningbox}[Shapes vs labelled trees]
$C_n$ counts \emph{shapes} (labels ignored); $n^{\,n-2}$ counts \emph{labelled} trees
(edges free, labels matter, not necessarily binary). Quoting the wrong one is a whole-
problem zero. Ask first: do node identities matter, and is binary structure required?
\end{warningbox}

\begin{warningbox}[Overflow and modulo]
Catalan and $n^{n-2}$ explode; in contests the modulus is part of the task. Never compute
$\binom{2n}{n}$ as raw integers -- use modular inverses or the DP, all under the modulus.
\end{warningbox}

\subsection{Solved Problems}

\solvedproblem{c01-unique-bst}{Unique Binary Search Trees}{LeetCode 96}
  {https://leetcode.com/problems/unique-binary-search-trees/}{LC Medium}
  {Catalan numbers, root-splitting DP}

\begin{problemstatement}[Unique BSTs]
Given $n$, count how many \emph{structurally unique} BSTs store the keys $1 \dots n$.
For $n = 3$ the answer is $5$.
\end{problemstatement}

\solstep{Understanding the Problem}
A BST's in-order traversal is the sorted keys, so once we pick which key is the root, the
smaller keys are forced left and the larger are forced right. Different shapes give
different trees, and two trees with the same shape but different keys can't both be valid
BSTs on $1..n$. So counting BSTs is exactly counting binary shapes on $n$ nodes -- the
Catalan number $C_n$ (\thmref{c01-catalan-shapes}).

\solstep{Key Observations}
\begin{observation}
Choosing root key $r$ splits keys into $r-1$ smaller and $n-r$ larger; the two sides form
independent subtrees.
\end{observation}
\begin{observation}
So $G(n) = \sum_{r=1}^{n} G(r-1)\,G(n-r)$ with $G(0)=G(1)=1$ -- the Catalan recurrence.
\end{observation}

\solstep{Solution}
Bottom-up DP over subtree size, summing the products of left and right counts, exactly
\coderef{c01-catalan}.
\begin{figure}[H]
\centering
\begin{tikzpicture}[diagram]
  \node[troot] (r2) at (0,0) {2};
  \node[tleaf] (r1) at (-1.2,-1.3) {1};
  \node[tleaf] (r3) at (1.2,-1.3) {3};
  \draw[tedge] (r2) -- (r1); \draw[tedge] (r2) -- (r3);
  \node[annot] at (-3.3,-1.3) {root 2};
  \node[annot] at (3.3,-1.3) {1 left, 1 right};
\end{tikzpicture}
\caption{One of the five BSTs on $\{1,2,3\}$: root $2$ forces $1$ left and $3$ right.
Pick root $1$ and both $2,3$ go right (two shapes), root $3$ mirrors it -- $5$ in total.}\label{fig:c01-unique-bst}
\end{figure}

\solstep{Proof of Correctness}
\begin{proofbox}
By the BST in-order property, key $r$ as root uses exactly keys $<r$ on the left and $>r$
on the right, and the subtrees are independent valid BSTs, contributing
$G(r-1)G(n-r)$ trees. The cases partition all BSTs by their root, so $G(n)=\sum_r
G(r-1)G(n-r)$. The DP computes this with base $G(0)=G(1)=1$ and returns $G(n)=C_n$.
\end{proofbox}

\solstep{Implementation}
\begin{lstlisting}[caption={Unique Binary Search Trees (LeetCode 96)}, label={lst:c01-unique-bst}]
#include <bits/stdc++.h>
using namespace std;

class Solution {
public:
    int numTrees(int n) {
        vector<int> dp(n + 1, 0);
        dp[0] = 1; dp[1] = 1;
        for (int nodes = 2; nodes <= n; ++nodes)
            for (int root = 1; root <= nodes; ++root)
                dp[nodes] += dp[root - 1] * dp[nodes - root];
        return dp[n];
    }
};
\end{lstlisting}

\solstep{Complexity}
\begin{complexitybox}
Time $O(n^2)$, memory $O(n)$. For LC $n \le 19$ the value fits in \code{int}; on larger
inputs or with a modulus, do the same DP with \code{long long} and reduce. A closed-form
$C_n$ with modular inverses is $O(n)$.
\end{complexitybox}

\solstep{Pitfalls and Variants}
\begin{warningbox}[Counting vs generating]
\probref{c01-unique-bst} only \emph{counts}; \emph{Unique BST II} (LC 95) builds every
tree, so the same recurrence is applied to node objects instead of integers -- $O(C_n)$
output, and recursion, not a flat DP.
\end{warningbox}

\subsection{Practice Problems}

\begin{practiceset}
  \practice{Unique Binary Search Trees II}{LeetCode 95}
    {https://leetcode.com/problems/unique-binary-search-trees-ii/}{LC Medium}
    {same root split, but build and combine actual left/right subtree lists.}
  \practice{Different Ways to Add Parentheses}{LeetCode 241}
    {https://leetcode.com/problems/different-ways-to-add-parentheses/}{LC Medium}
    {each operator as a root splits the expression -- a Catalan-shaped interval DP. See \chref{c10}.}
  \practicehh{Count Binary Trees With Factors}{LeetCode 852}
    {https://leetcode.com/problems/count-binary-trees-with-factors/}{LC Medium}
    {sort values; for each root count factor pairs that multiply to it.}
    {DP: answer[v] = 2 * sum over proper factor pairs (a,b) of answer[a]*answer[b]; memoise with a hash map of value to count.}
\end{practiceset}
